(a) Work over an affine neighborhood $\operatorname{Spec}B$ of $y=\mathfrak p$ and an affine open $\operatorname{Spec}A$ in its inverse image. With $S=B\setminus\mathfrak p$, we have $A\otimes_Bk(y)\cong S^{-1}(A/\mathfrak pA)$. Its primes correspond exactly to primes $\mathfrak q\subseteq A$ containing $\mathfrak pA$ and disjoint from $S$, that is, to those contracting to $\mathfrak p$. Under this correspondence, a distinguished open defined by $a/s$ is the intersection of the fiber with $D(a)\subseteq\operatorname{Spec}A$. Thus the map is a homeomorphism onto the fiber with its induced topology. These local identifications glue.

(b) For $y=(s-a)$, we have $X_y=\operatorname{Spec}k[t]/(t^2-a)$. If $a\ne0$ and $\operatorname{char}k\ne2$, then $t^2-a=(t-\sqrt a)(t+\sqrt a)$ has distinct factors, so the fiber is $\operatorname{Spec}(k\times k)$. At $a=0$, the fiber is the nonreduced one-point scheme $\operatorname{Spec}k[t]/(t^2)$. In characteristic $2$, $t^2-a=(t-\sqrt a)^2$, so every closed fiber is a nonreduced one-point scheme.

For the generic point $\eta$, we have $X_\eta=\operatorname{Spec}k(s)[t]/(t^2-s)=\operatorname{Spec}k(t)$, where $s$ acts as $t^2$. The polynomial $t^2-s$ is irreducible over $k(s)$ by Eisenstein at $s$, so $[k(t):k(s)]=2$.
